% Generated draft: requires mathematical and source review.
\documentclass[11pt]{article}
\usepackage[margin=25mm]{geometry}
\usepackage{fontspec}
\setmainfont{Latin Modern Roman}
\usepackage{amsmath,amssymb,mathtools,unicode-math}
\setmathfont{Latin Modern Math}
\usepackage{enumitem,needspace,xurl,hyperref}
\setlength{\parindent}{0pt}
\setlength{\parskip}{5pt}
\setlength{\emergencystretch}{3em}
\newcommand{\N}{\mathbb{N}}
\newcommand{\Z}{\mathbb{Z}}
\newcommand{\Q}{\mathbb{Q}}
\newcommand{\R}{\mathbb{R}}
\newcommand{\C}{\mathbb{C}}
\newcommand{\F}{\mathbb{F}}
\newcommand{\A}{\mathbb{A}}
\newcommand{\PP}{\mathbb{P}}
\newcommand{\E}{\mathbb{E}}
\newcommand{\eps}{\varepsilon}
\newcommand{\dd}{\,\mathrm{d}}
\begin{document}
\section*{UnsolvedMath ID: 800249}
\textbf{Source problem code:} catalogue entry
\subsection{Definitions and mathematical statement}
Use the standard interpretation of a complete hyperbolic manifold as a connected complex manifold whose Kobayashi distance is a genuine complete metric. The source does not explicitly define this terminology or prescribe a complex dimension.

Let \(\Delta=\{z\in\C:|z|<1\}\), equipped with the Poincar\'e distance
\[
\rho_\Delta(a,b)=\frac12\log\frac{1+\left|\frac{a-b}{1-\overline b a}\right|}{1-\left|\frac{a-b}{1-\overline b a}\right|}.
\]
For complex manifolds \(Y,Z\), write \(\operatorname{Hol}(Y,Z)\) for the set of holomorphic maps from \(Y\) to \(Z\). For a connected complex manifold \(X\), its Carath\'eodory pseudodistance is
\[
c_X(p,q)=\sup_{f\in\operatorname{Hol}(X,\Delta)}\rho_\Delta(f(p),f(q)).
\]
Its Kobayashi pseudodistance is
\[
k_X(p,q)=\inf\left\{\sum_{j=1}^{m}\rho_\Delta(a_j,b_j):
\begin{array}{l}
m\geq1,\quad a_j,b_j\in\Delta,\quad f_j\in\operatorname{Hol}(\Delta,X),\\
f_1(a_1)=p,\quad f_m(b_m)=q,\\
f_j(b_j)=f_{j+1}(a_{j+1})\quad(1\leq j<m)
\end{array}\right\}.
\]
Completeness means that every \(k_X\)-Cauchy sequence converges in the metric \(k_X\). A convex domain in \(\C^n\) means a nonempty open convex subset of \(\C^n\), with convexity taken in the underlying real vector space.

\textbf{Problem (M. Abate; Problem 48 in the supplied 1995 list).} Determine whether there exist an integer \(n\geq1\) and a connected complex manifold \(X\) of complex dimension \(n\) such that
\begin{enumerate}
\item \(k_X\) is a genuine complete metric;
\item \(k_X(p,q)=c_X(p,q)\) for every \(p,q\in X\);
\item there is no convex domain \(D\subset\C^n\) biholomorphic to \(X\).
\end{enumerate}
Construct such a manifold or prove that none exists. The conventions above make explicit the standard reading of the source; the source itself supplies no further hypotheses.
\subsection{Short English statement}
Does there exist a complete Kobayashi hyperbolic complex manifold on which the Kobayashi and Carath\'eodory distances agree for every pair of points, but which is not biholomorphic to any convex domain? The dimension is unrestricted. This is a statement-only draft requiring human mathematical review.
\subsection{Sources}
\begin{itemize}
\item Supplied local UnsolvedMath registry record, permanent ID \(800249\), exact source problem code \texttt{catalogue entry}, titled \emph{Several Complex Variables Problem List --- Problem 48}. The supplied registry labels the problem \emph{open}; that label is not an independently verified mathematical or literature assessment.
\item Buzzard, Fornaess, and Krantz, \emph{The Several Complex Variables Problem List} (1995), Problem 48. Supplied source URL: \url{https://arxiv.org/abs/math/9509202}. The record identifies the extraction as \texttt{arxiv\_latex} and reports access on 2026-10-16; these are supplied provenance details, not a claim of independent source access.
\item The supplied statement attributes the problem to M. Abate and gives the contact address \texttt{abate@vaxsns.infn.it}. The registry's proposer field names Buzzard, Fornaess, and Krantz, with proposed year 1995; their list attribution is retained separately from the statement's attribution to Abate.
\end{itemize}
\end{document}
